\section{Low-density cones and attainment}\label{sec:cones}

We now obtain the compactness needed to minimize among counterexamples.
The regularity argument has two induction parameters: the hypersurface
dimension $m$ in Theorem~\ref{thm:volume}, and the dimension of a Euclidean
cone.  Keeping them separate is essential.  The theorem in dimensions
less than $m$ will classify cones of dimension at most $m$, and these
classifications will give a smooth link for a cone of dimension $m+1$.

For a varifold $V$, its weight measure is
$\|V\|$ and its total mass, when finite, is $\mathbf M(V)$.
For a smooth embedded submanifold $S$, write $\mathbf v(S)$ for its
associated varifold with multiplicity one.
For an integral $k$-varifold in Euclidean space, $[V]$ denotes the
locally rectifiable flat chain with coefficients in $\mathbb Z_2$
obtained by reducing its integer multiplicity modulo two.  We call
$V$ \emph{cyclic modulo two} if $\partial[V]=0$ locally in the ambient
open set.  This convention applies to noncompact cones and their
translates, using locally flat chains as in
\cite[\S2.1 and Appendix]{White2009}.

\subsection{Coning, density, and the radial factor}

\begin{lemma}[The cone operation]\label{lem:cone-operation}
Let $d\geq1$ and let $V$ be a stationary integral $d$-varifold in the
unit sphere $S^{d+1}$.  Its Euclidean cone $\mathcal C(V)$ is a
stationary integral $(d+1)$-varifold in $\mathbb R^{d+2}$, and
\begin{equation}\label{eq:cones-mass}
 \|\mathcal C(V)\|(B_R(0))
   =\frac{R^{d+1}}{d+1}\mathbf M(V),
 \qquad
 \Theta(\mathcal C(V),0)=\frac{\mathbf M(V)}{\sigma_d}.
\end{equation}
If $V$ is a finite sum of smooth closed embedded minimal
hypersurfaces with positive integer multiplicities, then
$\mathcal C(V)$ is cyclic modulo two.
\end{lemma}

\begin{proof}
At a spherical point $x$ with tangent plane $P\subset x^\perp$, assign
to $rx$ the plane $\mathbb Rx\oplus P$ and the measure
$r^d\,dr\,dV(x,P)$.  This defines the cone and gives
\eqref{eq:cones-mass}.  The polar map is a diffeomorphism away from
the origin; rectifiability and integer multiplicity there follow
from the area formula.  The origin has zero mass.

For stationarity, first take a smooth ambient field $X$ supported
away from the origin and write
$X(rx)=f(r,x)x+Y_r(x)$, where $Y_r(x)\perp x$.
The trace of its derivative on the cone plane is
\[
 \operatorname{div}_{\mathbb Rx\oplus P}X(rx)
 =\partial_r f(r,x)+\frac{d}{r}f(r,x)
       +\frac1r\operatorname{div}^{S}_{P}Y_r(x).
\]
After multiplication by $r^d$ and integration, the first two terms
integrate to the endpoint values of $r^df$, and the last term
vanishes by spherical stationarity for each $r$.  Thus the cone is
stationary off the origin.  For a general compactly supported $X$,
insert a radial cutoff that is zero on $B_\varepsilon$ and one
outside $B_{2\varepsilon}$, with derivative bounded by
$c/\varepsilon$.  Formula~\eqref{eq:cones-mass} bounds the omitted
first-variation terms by
$c_X(\varepsilon^{d+1}+\varepsilon^d)$, which tends to zero.

For the last assertion, $[\mathcal C(V)]$ has zero boundary on
every annulus: an odd component multiplicity gives a smooth
mod-two product cylinder, and an even multiplicity gives zero.
Truncate this chain inside radius $\varepsilon$.  Its difference
from the full chain has mass at most
$\mathbf M(V)\varepsilon^{d+1}/(d+1)$, whereas its inner boundary
has mass at most $\mathbf M(V)\varepsilon^d$.
The truncated chains therefore converge locally flat to
$[\mathcal C(V)]$ and their boundaries converge locally flat to zero.
Continuity of the boundary operator proves
$\partial[\mathcal C(V)]=0$ also at the vertex.
\end{proof}

We use the stationary monotonicity formula, tangent-cone existence,
and integral-varifold closure of Allard
\cite[\S5 and Theorem~6.4]{Allard1972}; see also the formulations in
\cite[Chapter~4, \S3]{SimonGMT} and
\cite[Theorem~3.2]{White2009}.
In particular, on a fixed open set, locally bounded mass and locally
bounded first variation give a subsequence with an integral-varifold
limit.  For the stationary sequences below the first variation is
identically zero.  White's convergence theorem then gives the
additional implication
\begin{equation}\label{eq:cones-white}
 V_i\longrightarrow V,\quad \delta V_i=0,\quad
 \partial[V_i]=0
 \quad\Longrightarrow\quad \partial[V]=0,
\end{equation}
provided the masses are locally bounded.  Indeed, Allard gives
integrality of $V$; the first variations satisfy White's local bound
with bound zero; and the boundary chains converge to zero.
These are precisely the hypotheses of
\cite[Definition~3.1 and Theorem~3.3]{White2009}.
The same statement applies after any translations and dilations,
and hence to the blow-ups used below.

\begin{lemma}[Density at every support point]\label{lem:cones-density}
Let $V$ be a stationary integral $k$-varifold in an open set
$U\subset\mathbb R^n$.  At every $z\in\spt\|V\|$,
$\Theta(V,z)\geq1$.  If $V=\mathcal C$ is a cone in all of
$\mathbb R^n$, then
\begin{equation}\label{eq:cones-density-comparison}
 1\leq\Theta(\mathcal C,z)\leq\Theta(\mathcal C,0)
 \qquad (z\in\spt\|\mathcal C\|).
\end{equation}
Every tangent cone at such a point is nonzero and has density
$\Theta(\mathcal C,z)$ at its vertex.
\end{lemma}

\begin{proof}
Integrality gives density at least one at $\|V\|$-almost every point.
Every neighborhood of $z$ has positive mass, so choose such points
$x_i\to z$.  If $\overline{B_r(z)}\subset U$ and
$d_i=|x_i-z|<r$, monotonicity at $x_i$ and ball inclusion give
\[
 \|V\|(B_r(z))\geq\|V\|(B_{r-d_i}(x_i))
       \geq\omega_k(r-d_i)^k.
\]
Let $i\to\infty$, then $r\downarrow0$.  This proves the asserted
lower bound also at support points that are exceptional for
rectifiability.

For a cone, its homogeneous mass measure gives
$\|\mathcal C\|(B_R(0))=\omega_k\Theta(\mathcal C,0)R^k$.
For $R>|z|$, the inclusions
\[
 B_{R-|z|}(0)\subset B_R(z)\subset B_{R+|z|}(0)
\]
identify the limit of the mass ratios centered at $z$ as
$R\to\infty$ with $\Theta(\mathcal C,0)$.
Monotonicity proves the upper bound in
\eqref{eq:cones-density-comparison}.  Tangent-cone existence and
the monotonicity formula give the stated tangent density, which is
positive by the first part.
\end{proof}

\begin{lemma}[Splitting the radial line]\label{lem:cones-splitting}
Let $k\geq2$ and let $\mathcal C$ be a stationary integral
$k$-dimensional cone in $\mathbb R^{k+1}$.
For $0\ne z\in\spt\|\mathcal C\|$, every tangent cone $T$ at $z$
has the form
\[
 T=\mathbf v(\mathbb Re)\times W,\qquad e=z/|z|,
\]
where $W$ is a nonzero stationary integral $(k-1)$-dimensional
cone in $e^\perp\cong\mathbb R^k$ and
\begin{equation}\label{eq:cones-factor-density}
 \Theta(W,0)=\Theta(T,0)=\Theta(\mathcal C,z).
\end{equation}
If $\mathcal C$ is cyclic modulo two, so are $T$ and $W$.
\end{lemma}

\begin{proof}
Write $T_i=(\eta_{z,\rho_i})_\#\mathcal C\to T$, where
$\eta_{z,\rho_i}(y)=(y-z)/\rho_i$ and $\rho_i\downarrow0$.
For fixed $b\in\mathbb R$, dilation invariance of $\mathcal C$
by the positive factors $1+b\rho_i$ implies invariance of $T_i$
under
\[
 y\longmapsto(1+b\rho_i)y+bz.
\]
Passing to the limit shows that $T$ is invariant under translation
by $bz$.  Its weight consequently disintegrates as
$\|T\|=\mathcal L^1\otimes\nu$ on
$\mathbb Re\oplus e^\perp$.

We check that this is an integral varifold product.
At almost every rectifiable point of $T$, a measure blow-up is an
integer multiple of the measure on its approximate tangent plane.
These blow-ups retain invariance under translations in the
$e$-direction, so that plane contains $e$.
The tangential Jacobian of the coordinate
$\pi(te+w)=t$ is therefore one almost everywhere.
Rectifiable coarea gives, for almost every $t$, a
$(k-1)$-rectifiable slice with the original integer multiplicities;
its tangent plane is the intersection of the tangent plane of $T$
with $e^\perp$.  Product disintegration identifies the weight of
this slice, after translation to $e^\perp$, with $\nu$.
Thus $\nu$ defines an integral varifold $W$ and the tangent-plane
identification gives $T=\mathbf v(\mathbb Re)\times W$ as varifolds.
Here we use rectifiable coarea
\cite[\S3, p.~300]{Federer1978}.

For a compactly supported field $Y$ on $e^\perp$ and a compactly
supported smooth function $\chi$ on $\mathbb R$ with
$\int\chi\ne0$, test stationarity of $T$ with
$X(te+w)=\chi(t)Y(w)$.  Since $Y$ is perpendicular to $e$,
\[
 0=\delta T(X)=\left(\int_{\mathbb R}\chi(t)\,dt\right)\delta W(Y).
\]
Hence $W$ is stationary.  Dilation invariance of the tangent cone
$T$ and the product formula imply dilation invariance of $W$.
Writing $\theta=\Theta(W,0)$, integration over a ball gives
\[
 \|T\|(B_1)
 =\theta\omega_{k-1}
       \int_{-1}^1(1-t^2)^{(k-1)/2}\,dt
 =\theta\omega_k.
\]
This proves \eqref{eq:cones-factor-density};
Lemma~\ref{lem:cones-density} makes $W$ nonzero.

If $\partial[\mathcal C]=0$, every $T_i$ is cyclic, and
\eqref{eq:cones-white} implies $\partial[T]=0$.
The mod-two slicing identity for the boundary
\cite[\S3, equation~(1) and Proposition~3.1]{White1999} gives
\[
 \partial\langle[T],\pi,t\rangle
       =\langle\partial[T],\pi,t\rangle=0
 \quad\text{for almost every }t,
\]
where signs disappear modulo two. For the locally flat chain, apply
the identity on a countable exhaustion by relatively compact regions;
restriction compatibility gives one full-measure set of levels where
it holds locally everywhere. Choose one such level where the preceding
rectifiable slice description also holds. The slice there is the
translate of $[W]$ to $te+e^\perp$,
because $T$ is the product just established.  Translation back
proves $\partial[W]=0$.  No assertion about whether zero is a good
slicing level is needed.
\end{proof}

For clarity, we also record the precise regularity consequence
used in the dimension reduction.

\begin{lemma}[A multiplicity-one planar tangent]\label{lem:cones-allard}
Let $V$ be a stationary integral $k$-varifold, $k\geq2$, in an open
subset of Euclidean space.  If one tangent cone at $z$ is a
multiplicity-one plane, then in a neighborhood of $z$, $V$ is a
smooth embedded minimal submanifold with multiplicity one.
\end{lemma}

\begin{proof}
Apply Allard's interior regularity theorem
\cite[\S8]{Allard1972}, in the formulation of
\cite[Chapter~5, hypotheses~(5.1) and Theorem~5.2]{SimonGMT},
on a sufficiently small rescaled ball from
the sequence giving that tangent.  Its hypotheses hold as follows.
Integrality supplies the almost-everywhere density lower bound one.
Stationarity supplies zero generalized mean curvature and zero
first variation, hence the required mean-curvature bound for any
exponent greater than $k$.  Convergence to the multiplicity-one
plane supplies normalized mass arbitrarily close to $\omega_k$
and arbitrarily small height and tangent-plane excess on a fixed
smaller ball.  The center remains in the support.  Allard's theorem
therefore gives a single $C^{1,\alpha}$ graph near the center.
On that connected graph the stationary integral multiplicity is
constant; its tangent multiplicity is one, so the constant is one.
The higher-regularity conclusion of
\cite[Chapter~5, Theorem~6.1]{SimonGMT}, applied with the smoothly
prescribed mean curvature identically zero, makes the graph smooth;
that statement permits arbitrary codimension. This theorem has
no stability hypothesis and no upper bound on $k$.
\end{proof}

\Needspace{17\baselineskip}
\subsection{The cone induction}

\begin{proposition}[Low-density cones]\label{prop:low-density-cones}
Fix $m\geq2$ and assume Theorem~\ref{thm:volume} in dimensions
$2,\ldots,m-1$, with an empty assumption when $m=2$.
Let $\mathcal C$ be a nonzero stationary integral
$k$-dimensional cone in $\mathbb R^{k+1}$ such that
\begin{equation}\label{eq:cones-threshold}
 \Theta(\mathcal C,0)<\vartheta_m.
\end{equation}
When $m=2$ or $m=3$, also assume that $\mathcal C$ is cyclic
modulo two.  Then:
\begin{enumerate}[label=\textup{(\roman*)}]
 \item if $1\leq k\leq m$, $\mathcal C$ is a
 multiplicity-one hyperplane;
 \item if $k=m+1$, $\mathcal C$ is smooth embedded with
 multiplicity one away from its vertex.  Its link in $S^{m+1}$
 is a smooth closed connected embedded minimal hypersurface.
\end{enumerate}
\end{proposition}

\begin{proof}
Hold $m$ fixed throughout.  We prove (i) by induction in $k$ and
then obtain (ii) from the final induction step.

An integral one-dimensional cone is a finite union of rays with
positive integer weights: radial slicing of its conical
rectifiable measure gives a zero-dimensional integral measure on
the unit circle, and local finiteness makes its total weight finite.
Write the distinct unit ray directions as $u_j$, their weights
as $q_j$, and $N=\sum_jq_j$.  Stationarity at the origin says
\[
 \sum_jq_ju_j=0,\qquad \Theta(\mathcal C,0)=N/2.
\]
If $m\geq4$, Lemma~\ref{lem:clifford-ratios} gives
$\vartheta_m\leq3/2$, so the strict bound
\eqref{eq:cones-threshold} gives $N<3$.
Nonzero stationarity excludes $N=1$, leaving two opposite rays of
weight one.  If $m=2$ or $3$, the same lemma gives
$\vartheta_m<2$, hence $N<4$; cyclicity says that the boundary
coefficient $N$ at the origin is even.  Again $N=2$ and balance
gives a multiplicity-one line.  This proves the base case in both
branches.

Suppose next that $2\leq k\leq m+1$ and that (i) has been proved
in dimension $k-1$.  At any nonzero support point $z$,
Lemma~\ref{lem:cones-splitting} expresses a tangent as
$\mathbf v(\mathbb R)\times W$.  Its factor is nonzero, stationary and
integral, is cyclic when required, and has density at most
$\Theta(\mathcal C,0)<\vartheta_m$ by
Lemma~\ref{lem:cones-density}.  The induction hypothesis makes
$W$ a multiplicity-one plane.  Thus this tangent of $\mathcal C$
is a multiplicity-one plane, and
Lemma~\ref{lem:cones-allard} proves smooth embedded
multiplicity-one structure near $z$.

It follows that $\Lambda=\spt\|\mathcal C\|\cap S^k$ is a smooth
closed embedded $(k-1)$-manifold.  The radial direction lies in
the cone tangent plane and is normal to the sphere, so this
intersection is transverse.  Polar coarea gives
$\mathcal C=\mathcal C(\mathbf v(\Lambda))$ with multiplicity one.
Reversing the angular first-variation calculation in
Lemma~\ref{lem:cone-operation} shows that $\Lambda$ is minimal
in $S^k$.

Each connected component $\Lambda_j$ is compact and minimal,
and its cone is stationary.  Lemma~\ref{lem:cones-density}
applied to that cone, together with \eqref{eq:cones-mass}, yields
\begin{equation}\label{eq:cones-component}
 \frac{\Vol(\Lambda_j)}{\sigma_{k-1}}\geq1.
\end{equation}
Since the total cone density is less than $\vartheta_m<2$,
there is only one component.

If $k=2$, a closed connected embedded minimal curve in $S^2$ is a
great circle by the geodesic equation, so its cone is a plane.
If $3\leq k\leq m$ and $\Lambda$ were not totally geodesic, the
outer induction in dimension $k-1$ would give
\[
 \Theta(\mathcal C,0)
   =\frac{\Vol(\Lambda)}{\sigma_{k-1}}
   \geq\vartheta_{k-1}\geq\vartheta_m,
\]
contrary to \eqref{eq:cones-threshold}.  Thus $\Lambda$ is a
great sphere and $\mathcal C$ a multiplicity-one plane, completing
(i).  For $k=m+1$, the already proved plane assertion in dimension
$m$ gives exactly the smoothness and connectedness asserted in
(ii).  We do not apply the dimension-$m$ volume statement to this
last link.
\end{proof}

\begin{corollary}[Spherical stationary output]\label{cor:cones-spherical-output}
Under the outer induction hypothesis of
Proposition~\ref{prop:low-density-cones}, let $V\ne0$ be a
stationary integral $m$-varifold in $S^{m+1}$ with
$\mathbf M(V)<a_m$.  If $m=2$ or $3$, assume additionally that
$\mathcal C(V)$ is cyclic modulo two; this holds, in particular,
if $V$ has smooth closed embedded support with integer component
multiplicities.  Then $V=\mathbf v(\Sigma)$ for a smooth closed connected
embedded minimal hypersurface $\Sigma$.
If $\mathbf M(V)>\sigma_m$, this hypersurface is not an equator.
\end{corollary}

\begin{proof}
Lemma~\ref{lem:cone-operation} and the mass bound put
$\mathcal C(V)$ under part (ii) of
Proposition~\ref{prop:low-density-cones}.  Its smooth
multiplicity-one link identifies $V$ by the polar cone formula.
An equator has area $\sigma_m$.
\end{proof}

\subsection{Exact attainment}

\begin{proposition}[A least counterexample]\label{prop:attainment}
Fix $m\geq2$ and assume Theorem~\ref{thm:volume} in dimensions
$2,\ldots,m-1$.  Suppose the class $\mathcal E_m$ of smooth
closed connected embedded minimal hypersurfaces
$\Sigma\subset S^{m+1}$ satisfying
\[
 \Sigma\text{ is not totally geodesic},\qquad \Vol(\Sigma)<a_m
\]
is nonempty.  Then
\[
 H=\inf_{\Sigma\in\mathcal E_m}\Vol(\Sigma)
\]
is attained by some $\Sigma_*\in\mathcal E_m$, and
\begin{equation}\label{eq:cones-attainment-gap}
 \sigma_m<\gamma_m\leq H<a_m.
\end{equation}
Here $\gamma_m$ is the detection gap from
Lemma~\ref{lem:detection-gap}.
\end{proposition}

\begin{proof}
For every member of $\mathcal E_m$, the family of
Proposition~\ref{prop:lorentz-family} is detecting and has maximum
mass at most its area.  Lemma~\ref{lem:detection-gap} therefore
gives $\Vol(\Sigma)\geq\gamma_m>\sigma_m$.
Nonemptiness gives $H<a_m$, proving
\eqref{eq:cones-attainment-gap}.

Choose $\Sigma_i\in\mathcal E_m$ with areas $H_i\to H$ and set
$\mathcal C_i=\mathcal C(\mathbf v(\Sigma_i))$.
These cones are stationary, integral, and cyclic by
Lemma~\ref{lem:cone-operation}.  Their exact masses satisfy
\begin{equation}\label{eq:cones-exact-sequence}
 \|\mathcal C_i\|(B_R(0))=\frac{H_iR^{m+1}}{m+1}
 \qquad(R>0).
\end{equation}
Consequently their masses are uniformly bounded on every compact
set.  Allard's integral-varifold compactness theorem and a
diagonal subsequence give a stationary integral limit
$\mathcal C$ in all of $\mathbb R^{m+2}$.
Every fixed dilation leaves each $\mathcal C_i$ invariant, so by
continuity of pushforward it leaves $\mathcal C$ invariant as well.
White's theorem in the form \eqref{eq:cones-white} gives
$\partial[\mathcal C]=0$.

We verify that the limiting mass is exactly $H$ at the spherical
level.  The weight of the conical limit is homogeneous of degree
$m+1$.  It has zero mass on every sphere of positive radius:
otherwise dilation would produce infinitely many disjoint spheres
with uniformly positive mass in a fixed compact annulus.
Weak convergence can thus be evaluated on every centered ball.
Passing to the limit in \eqref{eq:cones-exact-sequence} yields
\[
 \|\mathcal C\|(B_R(0))=\frac{HR^{m+1}}{m+1},
 \qquad
 \Theta(\mathcal C,0)=\frac H{\sigma_m}<\vartheta_m.
\]
In particular $\mathcal C$ is nonzero.  This exact identity
retains the infimum value; a lower-semicontinuity bound would not
give the same conclusion by itself.

Proposition~\ref{prop:low-density-cones}(ii) now makes the link
$\Sigma_*=\spt\|\mathcal C\|\cap S^{m+1}$ a smooth closed
connected embedded minimal hypersurface, with multiplicity one.
Its area is exactly $H$.  The bounds
$\sigma_m<H<a_m$ exclude an equator and put $\Sigma_*$ in
$\mathcal E_m$, as required.
\end{proof}
